% Proposed insertion after Theorem~\ref{thm:rank} in paper_merged.tex
\begin{remark}[When are the assumptions approximately satisfied?]
\label{rem:rank_assumptions_empirical}
Theorem~\ref{thm:rank} is a local structural statement rather than a universal claim. Assumption~\ref{ass:align} is favored when the preconditioner tracks the dominant curvature directions of the layer and these directions vary slowly over training. This is approximately the regime created by BatchNorm in deep networks: BN suppresses Hessian outliers and reduces basis twisting, making the left singular directions of $G$ more stable relative to the eigenvectors of $P$. In architectures without BN, SDP plays a similar role at task boundaries by compressing the singular-value spectrum of the weights and empirically resetting the Hessian condition number, which makes both spectral alignment and order preservation more plausible on the next task. For Fisher-based K-FAC, CWN complements this effect in non-BN MLPs by removing mean-dominated directions, stabilizing the effective Fisher geometry seen by the update.

The assumptions map differently across optimizer families. For Hessian-based methods such as AdaHessian, SophiaH, and SASSHA, they correspond to the approximate alignment between the dominant gradient subspace and the leading Hessian eigenspaces used for preconditioning. For Fisher-based methods such as K-FAC/NGD, they correspond to the preconditioned Fisher block acting as a near-whitening operator in a slowly varying feature basis; CWN is specifically designed to make this approximation more realistic when BN is absent. A detailed proof under Assumptions~\ref{ass:align}--\ref{ass:order} is given in Appendix~\ref{app:proof_rank}.
\end{remark}

% Proposed insertion near the beginning of Section~\ref{sec:experiments}
\paragraph{Experiments as assumption diagnostics.}
Beyond performance comparison, our experiments also probe when Assumptions~\ref{ass:align}--\ref{ass:order} are close to valid in practice. The deep BN setting (CIFAR-100 with ResNet-18) is the regime where the assumptions are expected to hold most closely, and indeed second-order methods already preserve rank well even without SDP. The no-BN settings reveal the opposite regime: when spectral conditioning deteriorates, the assumptions fail visibly through dormant-unit growth and stable-rank collapse. Adding SDP restores the relevant spectral regularity for Hessian-based methods, while CWN plays the analogous stabilizing role for Fisher-based K-FAC on non-BN MLPs. We therefore interpret the plasticity metrics, stable-rank trajectories, and the Hessian condition-number ablation as empirical signatures of how closely the assumptions behind Theorem~\ref{thm:rank} are satisfied.
